خیر، طراحی یک مدل یکپارچه برای مفهوم \lr{Customer} در کل سیستم یک اشتباه معماری است و در \lr{DDD} یک \lr{Anti-pattern} محسوب می‌شود. تلاش برای گنجاندن تمام ویژگی‌های مربوط به فروش، حسابداری، انبار و پشتیبانی در یک موجودیت واحد، منجر به ایجاد یک \lr{God Class} می‌شود. این رویکرد باعث درهم‌تنیدگی شدید تیم‌ها و کدها شده و نگهداری سیستم را به شدت دشوار می‌کند؛ زیرا هر تغییری از سوی یک دپارتمان می‌تواند بر کار سایر دپارتمان‌ها تأثیر منفی بگذارد و توافق بر سر یک زبان مشترک را عملاً غیرممکن می‌سازد.

بر اساس مفهوم \lr{Bounded Context} (زمینه محدود)، سیستم باید به مرزهای معنایی کاملاً مشخص و ایزوله تقسیم شود که در اینجا هر دپارتمان (\lr{Sales}، \lr{Accounting} و \lr{Support}) یک \lr{Bounded Context} مجزا به شمار می‌رود. در داخل هر یک از این مرزها، موجودیت \lr{Customer} معنای خاص خود را دارد و تنها شامل ویژگی‌هایی است که برای آن دپارتمان اهمیت دارد. به این ترتیب، ما در سطح کد چندین کلاس مجزا به نام \lr{Customer} (یا نام‌های دقیق‌تر مانند \lr{SalesCustomer} و \lr{AccountingCustomer}) خواهیم داشت که هر کدام با \lr{Ubiquitous Language} کانتکست خودشان طراحی شده‌اند و تیم‌های توسعه می‌توانند آن‌ها را کاملاً مستقل از یکدیگر تکامل دهند.

برای یکپارچه‌سازی و ارتباط بین این مدل‌های متفاوت در کانتکست‌های مختلف، باید از الگوهای \lr{Context Mapping} استفاده کرد. رایج‌ترین و بهترین رویکرد، اشتراک‌گذاری یک شناسه یکتای جهانی (\lr{Global Customer ID}) بین تمامی سیستم‌هاست. ارتباطات نیز غالباً از طریق معماری \lr{Event-Driven} پیاده‌سازی می‌شود؛ به این صورت که وقتی مشتری جدیدی ثبت‌نام می‌کند، یک \lr{Domain Event} مانند \lr{CustomerRegistered} منتشر می‌شود و سایر کانتکست‌ها (مثل حسابداری یا پشتیبانی) با دریافت این رویداد، مدل اختصاصی خودشان از آن مشتری را با همان شناسه مشترک در دیتابیس خود ایجاد می‌کنند. همچنین در مواقعی که یک کانتکست نیاز به خواندن مستقیم داده از کانتکست دیگری با مدلی متفاوت دارد، از الگوی \lr{Anti-Corruption Layer} استفاده می‌شود تا داده‌های دریافتی، پیش از ورود به مرز کانتکست مقصد، به مدل بومی آن ترجمه شوند.